\documentclass[11pt,a4paper]{article}
\usepackage[T1]{fontenc}
\usepackage[utf8]{inputenc}
\usepackage{lmodern,amsmath,amssymb}
\usepackage[margin=25mm]{geometry}
\usepackage{longtable,booktabs,array,calc}
\usepackage[hidelinks]{hyperref}
\hypersetup{pdftitle={A universal action floor needs a fixed action unit},pdfauthor={navstokgap research working notes}}
\newcommand{\tightlist}{\setlength{\itemsep}{0pt}\setlength{\parskip}{0pt}}
\setlength{\emergencystretch}{3em}
\setcounter{secnumdepth}{0}
\begin{document}
\hypertarget{a-universal-action-floor-needs-a-fixed-action-unit-and-no-admitted-similarity}{%
\section{A universal action floor needs a fixed action unit and no
admitted
similarity}\label{a-universal-action-floor-needs-a-fixed-action-unit-and-no-admitted-similarity}}

A positive action floor shared by all admitted masses and preparations
is a dimensionless multiple of a product of the model's fixed constants
with action units. A model whose fixed constants admit no such product
has floor zero or infinity, and a model whose admitted class is
preserved by a one-parameter similarity that rescales the observable has
floor zero even when such a product exists. These two elementary
criteria decide the recorded action-selection countertests from C002 to
Q13 before calculation, and they identify the single mass-independent
action unit available to classical electrodynamics:
\(k_e/c=e^2/(4\pi\epsilon_0c)\), which equals \(\alpha\hbar\) and is
exactly the excluded angular-action infimum of the relativistic Kepler
threshold C052. The selection question therefore separates into the
premise fixing a charge unit and the factor \(1/\alpha\) between the
electromagnetic and quantum action units. This is Q14's consolidation;
it promotes no ledger claim.

\hypertarget{setting}{%
\subsection{1. Setting}\label{setting}}

A model class is specified by fixed constants \(c_1,\dots,c_r\), a set
of admitted masses, a class of admitted preparations, and dynamics. Each
fixed constant has a dimension exponent vector \(d_i\in\mathbb R^3\)
with respect to mass, length and time. Charge enters only through
Coulomb energies \(q_iq_j/(4\pi\epsilon_0r)\), so
\(k=q_iq_j/(4\pi\epsilon_0)\) has dimension energy times length,
\(d_k=(1,3,-2)\), and no separate charge dimension is needed. The action
dimension vector is \(d_A=(1,2,-1)\). An action observable \(A\) assigns
a value in \([0,\infty]\) to each admitted solution with its
preparation. The attained set \(\mathcal A(c)\subset[0,\infty]\)
collects its values over all admitted masses, preparations and solutions
at constants \(c=(c_1,\dots,c_r)\), and the floor is

\[g(c)=\inf\{A\in\mathcal A(c):A>0\},\]

with \(g=\infty\) when no positive value is attained.

Two premises are used. \textbf{Dimensional homogeneity:} the admitted
class, dynamics and observable are defined by dimensionally homogeneous
relations. Consequently a change of units
\(\lambda=(\lambda_M,\lambda_L,\lambda_T) \in(0,\infty)^3\), written
\(\lambda^{d}=\lambda_M^{d_M}\lambda_L^{d_L} \lambda_T^{d_T}\), carries
the model at constants \(c\) to the model at constants
\(\lambda^{d_i}c_i\) and multiplies every attained action by
\(\lambda^{d_A}\):

\[\mathcal A(\lambda^{d_1}c_1,\dots,\lambda^{d_r}c_r)
 =\lambda^{d_A}\,\mathcal A(c).\tag{1}\]

\textbf{Universality:} the floor depends on the fixed constants only.
Admitted masses and preparation parameters are ranged over inside the
infimum, so they are absent from the argument list of \(g\). This is
obligation 4 of the \href{../research/ACTION_FIELD_TARGET.md}{action
target} written as a hypothesis. A fixed bath mass, cutoff, apparatus
length or supplied coefficient is a fixed constant when the model holds
it fixed for all admitted preparations.

\hypertarget{theorem-a-a-universal-floor-is-a-multiple-of-an-action-product}{%
\subsection{2. Theorem A: a universal floor is a multiple of an action
product}\label{theorem-a-a-universal-floor-is-a-multiple-of-an-action-product}}

\textbf{Theorem A.} Under dimensional homogeneity and universality:

\begin{enumerate}
\def\labelenumi{(\roman{enumi})}
\item
  If \(d_A\) lies outside the linear span of \(d_1,\dots,d_r\), then
  \(g(c)\in\{0,\infty\}\) for every \(c\).
\item
  If \(d_A=\sum_ia_id_i\), then \(\Pi(c)=\prod_ic_i^{a_i}\) has action
  units and \(g(c)=\Pi(c)\,F(\pi_1,\dots,\pi_s)\), where
  \(\pi_1,\dots,\pi_s\) is a basis of the dimensionless products of the
  constants and \(F\) takes values in \([0,\infty]\). When the \(d_i\)
  are linearly independent there is no dimensionless product and
  \(g=C\,\Pi\) for one number \(C\in[0,\infty]\).
\end{enumerate}

\emph{Proof.} Equation (1) and universality give
\(g(\lambda^{d_1}c_1,\dots,\lambda^{d_r}c_r)=\lambda^{d_A}g(c)\) for all
\(\lambda\). Write \(\lambda=e^{\ell}\) componentwise with
\(\ell\in\mathbb R^3\), so \(\lambda^{d}=e^{\ell\cdot d}\). For (i), the
orthogonal complement of \(\mathrm{span}\{d_i\}\) is not contained in
\(d_A^\perp\), because in finite dimension \(V^\perp\subset w^\perp\) is
equivalent to \(w\in V\). Choose \(\ell\) with \(\ell\cdot d_i=0\) for
all \(i\) and \(s=\ell\cdot d_A\neq0\). Then every constant is unchanged
and \(g(c)=e^{ts}g(c)\) for all real \(t\), which forces
\(g(c)\in\{0,\infty\}\). For (ii), \(\Pi\) has dimension vector
\(\sum_ia_id_i=d_A\), so \(g/\Pi\) is invariant under all unit changes;
an invariant function of the constants is a function of a basis of
dimensionless products, which is the Buckingham reduction. If the
\(d_i\) are independent the only dimensionless product is the constant
\(1\). \(\square\)

The theorem gives the form of a floor and its possible mass dependence,
and leaves the positivity of \(F\) to dynamics. Its content for
selection is the contrapositive of (i): a model that holds no
action-dimensional combination of its fixed constants cannot exclude
zero universally, whatever its dynamics.

\hypertarget{theorem-b-an-admitted-similarity-forces-floor-zero}{%
\subsection{3. Theorem B: an admitted similarity forces floor
zero}\label{theorem-b-an-admitted-similarity-forces-floor-zero}}

\textbf{Theorem B.} Fix the constants. Let \(\mathcal S\) be the
admitted set of solutions with their preparations and let
\(T_a:\mathcal S\to\mathcal S\), \(a\in(0,\infty)\), satisfy
\(A(T_as)=a^{\,\delta}A(s)\) for some \(\delta\neq0\). If some
\(s\in\mathcal S\) has \(0<A(s)<\infty\), then \(g=0\), and the attained
positive values are also unbounded above. If \(A\) takes one common
value \(g\) on \(\mathcal S\), then \(g\in\{0,\infty\}\).

\emph{Proof.} \(A(T_as)=a^{\delta}A(s)\) runs through \((0,\infty)\) as
\(a\) does. \(\square\)

The map \(T_a\) must preserve the admitted class at fixed constants. The
three standard instances are amplitude scaling of a linear dynamics with
a preparation class closed under multiplication by \(a\) (quadratic
actions, \(\delta=2\)), velocity scaling of a free or collisional
preparation (\(\delta=2\) for a plateau proportional to \(u^2\)), and
the space--time contraction \(q_a(t)=aq(t/a)\) of
\href{action-scale-dilation.md}{C056}, which is a similarity between
different external potentials (\(\delta=1\)); C057 is Theorem B for that
map with the class enlarged to be closed under it. A dynamics that is
linear near an equilibrium supplies \(T_a\) asymptotically: the
small-radius circles of C058 have \(\ell_R\sim\sqrt{mk_s}R^2\) because
the confining potential is harmonic at its minimum, and the harmonic
oscillator has the exact amplitude similarity.

Theorems A and B are independent necessary conditions. Q02's regulated
oscillator holds \(q\), \(m\), \(\epsilon_0\) and \(c\) fixed, so
\(q^2/(\epsilon_0c)\) is an available action product; its linear
response with a preparation class containing every field amplitude
admits \(E\mapsto aE\), \(x\mapsto ax\), \(J_*\mapsto a^2J_*\), and
Theorem B gives floor zero. Fixing the amplitude \(K\) instead removes
the similarity by adding an action-dimensional constant, which is what
``supplied coefficient'' means throughout the ledger.

\hypertarget{the-recorded-countertests-and-bounds-under-the-two-criteria}{%
\subsection{4. The recorded countertests and bounds under the two
criteria}\label{the-recorded-countertests-and-bounds-under-the-two-criteria}}

\begin{longtable}[]{@{}
  >{\raggedright\arraybackslash}p{(\columnwidth - 8\tabcolsep) * \real{0.2000}}
  >{\raggedright\arraybackslash}p{(\columnwidth - 8\tabcolsep) * \real{0.2000}}
  >{\raggedright\arraybackslash}p{(\columnwidth - 8\tabcolsep) * \real{0.2000}}
  >{\raggedright\arraybackslash}p{(\columnwidth - 8\tabcolsep) * \real{0.2000}}
  >{\raggedright\arraybackslash}p{(\columnwidth - 8\tabcolsep) * \real{0.2000}}@{}}
\toprule\noalign{}
\begin{minipage}[b]{\linewidth}\raggedright
Result
\end{minipage} & \begin{minipage}[b]{\linewidth}\raggedright
Fixed constants
\end{minipage} & \begin{minipage}[b]{\linewidth}\raggedright
Action product
\end{minipage} & \begin{minipage}[b]{\linewidth}\raggedright
Admitted similarity
\end{minipage} & \begin{minipage}[b]{\linewidth}\raggedright
Criterion
\end{minipage} \\
\midrule\noalign{}
\endhead
\bottomrule\noalign{}
\endlastfoot
C002 constant-force variations & \(m\), \(F\), \(T\): \((1,0,0)\),
\((1,1,-2)\), \((0,0,1)\) & \(F^2T^3/m\) & \(\eta\mapsto a\eta\),
\(\Delta S\mapsto a^2\Delta S\) & B \\
C024 collision plateau & \(m\), \(M\), \(\nu\) & none without a fixed
speed & incoming velocities \(\mapsto\epsilon u\),
\(H_*\mapsto\epsilon^2H_*\) & A and B \\
C036 composition & reference \(m_0\), \(u_0\), \(\lambda_0\) &
\(m_0u_0^2/\lambda_0\) & none once the reference is fixed & A:
\(g=C\,\Pi\) \\
C056--C057 contractions & \(m\), \(u\), no length & none &
\(q_a(t)=aq(t/a)\), \(\delta=1\) & A and B \\
C058 fixed force-bounded potential & \(m\), \(k_s\), \(b\), \(c\) &
\(\sqrt{mk_s}\,b^2\) & harmonic core, \(\ell_R\sim\sqrt{mk_s}R^2\) & B,
asymptotic \\
C060--C061 speed floor, force ceiling & \(P_*\), \(v_*\), \(F_{\max}\):
\((1,1,-1)\), \((0,1,-1)\), \((1,1,-2)\) & \(P_*^2v_*/F_{\max}\), unique
& none: the speed floor breaks amplitude scaling & A: \(g=C\,\Pi\),
\(C>0\) \\
C052 relativistic Kepler & \(k\), \(c\): \((1,3,-2)\), \((0,1,-1)\) &
\(k/c\), unique, mass-independent & none: \(r\mapsto\lambda r\),
\(p\mapsto p/\lambda\), \(m\mapsto m/\lambda\), \(t\mapsto\lambda t\)
fixes \(L\) & A: \(g=k/c\), \(C=1\) \\
C053 softened core & \(k\), \(c\), \(a\), \(m\) & \(k/c\) & none & A:
\(g=(k/c)F(amc^2/k)\), \(F\equiv0\) \\
Q02 radiation balance & \(q\), \(m\), \(\epsilon_0\), \(c\), \(w\),
\(\Lambda\) & \(q^2/(\epsilon_0c)\), \(mc^2/w\) & field amplitude,
\(\delta=2\) & B \\
Q04 topological sector & \(\rho\) (energy), \(c\): \((1,2,-2)\),
\((0,1,-1)\) & none & radius \(R\) of the degree-one family & A \\
Q08 string interference & impedance \(Z\), lengths, \(c\) & \(Z\ell^2\)
& pulse amplitude, \(\delta=2\) & B \\
Q13 CHSH witness & probabilities only & none & none needed & A \\
\end{longtable}

Each row uses the constants the cited note holds fixed; a row that lists
a supplied coefficient among the constants moves to the C036 pattern.
The remaining Q-series tests end in a free common multiplier (Q05--Q07)
or a dimensionless statistic (Q09--Q12) and fall under the same two
criteria. The G-track results are in the same position on the other
side: a quantum Hamiltonian with couplings in energy units and a lattice
spacing has no action product without a fixed time or a supplied
\(\hbar\), which is the recorded status of C126 and G05.

\hypertarget{which-classical-constants-supply-a-mass-independent-action-unit}{%
\subsection{5. Which classical constants supply a mass-independent
action
unit}\label{which-classical-constants-supply-a-mass-independent-action-unit}}

\textbf{Proposition.} Let the fixed constants be
\(k_e=e^2/(4\pi\epsilon_0)\), \(c\) and \(G\), with dimension vectors
\((1,3,-2)\), \((0,1,-1)\) and \((-1,3,-2)\). The unique product with
action units is \(k_e/c\). Without \(k_e\) there is none; admitting a
mass \(m\) adds \(Gm^2/c\) and its mass-dependent relatives.

\emph{Proof.} The three vectors are linearly independent, so the
exponents \((a,b,\gamma)\) with
\(a(1,3,-2)+b(0,1,-1)+\gamma(-1,3,-2)=(1,2,-1)\) are unique. Adding the
length and time components gives \(a+\gamma=1\); the mass component
gives \(a-\gamma=1\); hence \(a=1\), \(\gamma=0\), \(b=-1\). With \(G\)
and \(c\) alone the mass component forces \(\gamma=-1\), and then the
length and time components are inconsistent. With \(m\) admitted,
\(Gm^2/c\) has dimension \((-1+2,\,3-1,\,-2+1)=(1,2,-1)\). \(\square\)

In quantum units \(k_e/c=\alpha\hbar\) with
\(\alpha=e^2/(4\pi\epsilon_0\hbar c)\approx1/137\). Thus the fixed
constants of classical mechanics, electrodynamics and gravitation admit
exactly one mass-independent action, and it exists because the charge is
a fixed unit. Charge quantization is the classical premise that supplies
an action unit; the value it supplies is \(\alpha\) times the quantum
one. A fixed length together with a mass unit and \(c\) would supply
another, \(m_0c\,\ell_0\), at the price of a mass unit and a fundamental
length, both absent from the models in use.

\hypertarget{the-relativistic-kepler-threshold-through-the-five-obligations}{%
\subsection{6. The relativistic Kepler threshold through the five
obligations}\label{the-relativistic-kepler-threshold-through-the-five-obligations}}

C052 already realizes the unique electromagnetic unit. For one body in
the singular potential \(-k/r\) with \(k=k_e\) for two elementary
charges, regular bound orbits exist exactly for \(|L|>k/c\).

\begin{enumerate}
\def\labelenumi{\arabic{enumi}.}
\tightlist
\item
  \textbf{Definition and units.} The angular action \(|L|\) of regular
  bound orbits; \(k/c\) has action units.
\item
  \textbf{Exclusion of zero.} The excluded infimum \(k/c\) follows from
  the singular core, finite speed and binding (C052); a softened core
  removes it (C053). The premise excluding zero is the \(1/r\)
  singularity at fixed \(k\) and \(c\).
\item
  \textbf{Convergence.} No limit is required: the threshold is exact at
  fixed constants, and closes only as \(c\to\infty\) or
  \(k\downarrow0\).
\item
  \textbf{Universality.} The threshold is independent of the mass, as
  the explicit formula and the similarity in Section 4 show, and it
  bounds every regular bound preparation. It scales with the charge
  product: for a centre of charge \(Ze\) it is \(Z\alpha\hbar\). The
  model excludes recoil, radiation and a dynamical second body.
\item
  \textbf{Quantum role.} None is derived. The floor is \(\alpha\hbar\),
  so a factor \(1/\alpha\) separates it from the quantum unit, and no
  premise of the model supplies that factor.
\end{enumerate}

Remark, audited. The quantum Coulomb problems place their collapse
thresholds at \(Z\alpha=l+1/2\) (Klein--Gordon) and
\(Z\alpha=|\kappa|=j+1/2\) (Dirac). Both are the classical condition
\(|L|>k/c\) with \(|L|\) measured as \(\hbar(l+1/2)\) or
\(\hbar|\kappa|\), and C052's circular energies at \(|L|=\hbar|\kappa|\)
are exactly the Dirac levels with zero radial quantum number.
\href{../reviews/kepler-collapse-Q14.md}{The Q14 review} derives this
from C052's energy identity, whose \(1/r^2\) coefficient
\(c^2\ell^2-k^2\) is what the wave equations' indicial equations test,
and records the primary sources; it is the Sommerfeld--Dirac coincidence
of 1916 and 1928. The only quantum input is the unit of angular action
at the plunge boundary, so the factor \(1/\alpha\) of item 5 says that
the Dirac ground state reaches that boundary at \(Z=1/\alpha\). The
correspondence is ledger claim C130; Theorems A and B remain
exploratory.

\hypertarget{strategic-consequence}{%
\subsection{7. Strategic consequence}\label{strategic-consequence}}

Before calculating a selection candidate, list its fixed constants and
check the two criteria. If no action product exists, or an admitted
similarity rescales the observable, the candidate can supply at most a
spectral shape (Q02), a conditional bound (C060) or a dimensionless
witness (Q13), and the calculation may be skipped or bounded
accordingly.

The one classical route with a mass-independent positive floor is
electromagnetic: a fixed charge unit, relativity and a singular Coulomb
core give \(\alpha\hbar\). The open premises of the selection track are
therefore the origin of a fixed charge unit and the factor \(1/\alpha\);
both lie outside classical dynamics, and the second is the measured
fine-structure constant. A return to a physical selection mechanism
should start from a model whose constants are \(e\), \(\epsilon_0\),
\(c\) and masses, with a nonlinear coupling that breaks amplitude
similarity, for example the radiating two-body Coulomb problem that
\href{../reviews/radiation-noise-Q02.md}{Q02's review} names as its
remaining dependency; its radiation-reaction length \(k_e/(mc^2)\) and
action unit \(k_e/c\) are then both fixed. Linear reservoirs,
contraction-closed classes and dimensionless statistics remain parked by
the criteria above.

Proof status: elementary written proofs above; the dimensional reduction
is the established Buckingham theorem. Literature status: dimensional
analysis and mechanical similarity are established methods, and the
identification \(k_e/c=\alpha\hbar\) is standard; no novelty is
asserted. A bounded librarian comparison and written review are required
before any ledger promotion.
\end{document}
